\section{The energy gain assigned to defect seeds}
\label{sec:seed-energy}
The exponential gain in Section~\ref{sec:packet-estimates} comes from
a lower bound on the energy assigned to a finite set of defects. Fine
plaquette defects, predictor defects and coarse flags require different
arguments. In particular, some coarse terms can be negative. The proof
retains a reserve from the very large fine and predictor defects to pay
for these terms. We give the local algebra and the final allocation,
keeping the geometric and local comparison inputs explicit.

\subsection{A quartic bound for the curvature defect}
Use the real Frobenius inner product in the defining unitary matrix
representation of the group. Write $\mathfrak g$ for its tangent space
at the identity and $\Pi$ for orthogonal projection onto $\mathfrak g$.
Let $D>0$ be the colour normalization, so that
\[
 \|U-1\|=D\,r(U),\qquad
 \|\operatorname{col}(v)\|=D\|v\|.
\]
Here $r(U)$ is the group chord and $\operatorname{col}$ is the
isometric identification up to the factor $D$. The curvature coordinate
$F(U)$ is obtained by projecting $(U-U^*)/2$ onto $\mathfrak g$.
Since the Hermitian part of $U-1$ is orthogonal, for the real Frobenius
inner product, to every skew-Hermitian matrix,
\[
 \operatorname{col}(F(U))=\Pi(U-1).
\]
To check the orthogonality, if $H^*=H$ and $X^*=-X$, cyclicity of the
trace shows that $\operatorname{Re}\operatorname{tr}(HX^*)$ is its own
negative. Pythagoras therefore gives the exact identity
\begin{equation}\label{eq:seed-normal-defect}
 D^2\{r(U)^2-\|F(U)\|^2\}
       =\|(1-\Pi)(U-1)\|^2.
\end{equation}
In particular $\|F(U)\|\le r(U)$ for every group element.

The local tubular projection in the group data is analytic, equals
the identity on the group near $1$, and has derivative $\Pi$ at $1$.
Taylor's formula on a sufficiently small ambient ball consequently
gives
\[
 \|(1-\Pi)(U-1)\|\le C\|U-1\|^2
\]
for nearby group elements. Explicitly, bound the second derivative of
the projection on a smaller closed ball and integrate its second-order
Taylor remainder along the segment from $1$ to $U$. This proves the
inequality with one constant $C>0$. Combining it with
\eqref{eq:seed-normal-defect} gives constants $K,\delta>0$, depending
only on the group data, such that
\begin{equation}\label{eq:seed-curvature-bound}
 r^2-Kr^4\le\|F(U)\|^2\le r^2,
 \qquad \tfrac34r^2\le\|F(U)\|^2
 \quad\text{when }r=r(U)\le\delta.
\end{equation}
One can take $K=C^2D^2$ and decrease $\delta$ until $K\delta^2\le1/4$.
Thus the nonnegative defect
$\mathfrak f(U)=(r(U)^2-\|F(U)\|^2)/2$ is at most $Kr(U)^4/2$
in that neighbourhood.

This proves the group estimate needed for the coarse energy row.
Let $a_c>0$ be its quadratic coefficient and $\lambda_c>0$ its linear
coefficient, with $\lambda_c\le a_c/8$. At coarse threshold $s\ge0$,
the clipping function and the radial attenuation are
\[
 q_s(r)=\begin{cases}a_c r^2/2,&r\le s,\\\lambda_c sr,&r>s,\end{cases}
 \qquad
 \alpha_s(r)=\begin{cases}1,&r=0,\\\min\{1,s/r\},&r>0.\end{cases}
\]
Cancellation of the other coarse quadratic components gives the row
\begin{equation}\label{eq:seed-coarse-row}
 C_s(U)=a_c(\alpha_s-\alpha_s^2/2)\|F(U)\|^2-q_s(r(U)).
\end{equation}
For $0\le s\le\delta$ this has three useful bounds:
\begin{equation}\label{eq:seed-three-ranges}
 \begin{array}{ll}
 r\le s:& C_s(U)\ge-B_0s^4,\quad B_0=a_cK/2,\\
 s<r\le\delta:&C_s(U)\ge(a_c/4)sr,\\
 r>\delta:&C_s(U)\ge-\lambda_c sr.
 \end{array}
\end{equation}
For the first, $\alpha_s=1$ and use the quartic defect. In the other
two ranges $\alpha_s=s/r\in[0,1]$, so
$\alpha_s-\alpha_s^2/2\ge s/(2r)$. In the middle range this gives
$C_s\ge(3a_c/8-\lambda_c)sr\ge a_csr/4$; in the last range drop the
nonnegative curvature term. These arguments include $s=0$.

\subsection{The fine row and the discarded quadratic component}
Let the fine threshold be $t>0$. Write $a_f,\lambda_f>0$ for the fine
quadratic and linear coefficients, and let $r_i$ be the fine plaquette
chords. The direct part of its energy row is
\[
 D_i=q_t(r_i)+\tfrac12\|v_i\|^2-\langle v_i,y_i\rangle,
 \qquad \|y_i\|\le\sqrt{a_f}\,r_i,
\]
where $q_t$ uses coefficients $a_f,\lambda_f$. The second part is
\[
 H_i=\mathbf1_{\{i\text{ retained}\}}\tfrac12\|u_i\|^2
       +\tfrac12\|v'_i\|^2-\langle v'_i,u_i\rangle.
\]
Their sum is the fine row. The guide bounds are
$\|v_i\|,\|v'_i\|\le\epsilon t$, with
$\epsilon=C/L^2$. The guide constant $C\ge0$, the coefficients $a_f,\lambda_f>0$,
and the fine clipping radius $R>0$ are fixed before $L$ is chosen.

Let $\mathcal F$ be the listed fine defects and
$\mathcal S=\{i\in\mathcal F:r_i>t\}$ their large part. The local
sector conditions ensure $r_i\ge t/4$ on $\mathcal F$ and
$r_i\le t/2$ on unlisted plaquettes in the relevant neighbourhood.
For $r_i\le t$, Cauchy--Schwarz gives
\[
 D_i\ge\tfrac12(\sqrt{a_f}\,r_i-\|v_i\|)^2\ge0.
\]
If additionally $i\in\mathcal F\setminus\mathcal S$ and
$\epsilon\le\sqrt{a_f}/8$, the last square is at least
$a_ft^2/128$. For $i\in\mathcal S$, the linear part of the clipping
and $\epsilon\sqrt{a_f}\le\lambda_f/4$ give
$D_i\ge(3\lambda_f/4)tr_i$. Summing over the fine rows in the
$2B$ neighbourhood of the seeds therefore gives
\begin{equation}\label{eq:seed-direct-income}
 \sum_i D_i\ge \frac{3\lambda_f}{4}t\sum_{i\in\mathcal S}r_i
             +\frac{a_f}{128}t^2|\mathcal F\setminus\mathcal S|.
\end{equation}

The second component is nonnegative on retained rows, since there it
is $\|u_i-v'_i\|^2/2$. It is not necessarily nonnegative on discarded
rows. We count these rows on the original fine lattice. A discarded
row at $i$ has a witness $j$ with
\[
 t\exp\{d(i,j)/R\}<r_j.
\]
The finite-range locality conditions ensure that each such witness
belongs to $\mathcal S$. For $y\ge1$, the number of fine plaquette
addresses within distance $R\log y$ of a point is at most
\[
 C_R y,\qquad C_R=(8R+1)^4\#\{\text{plaquette axes}\}.
\]
Indeed $1+(\log y)/4\le y^{1/4}$ implies
$2R\log y+1\le(8R+1)y^{1/4}$, and the lattice ball has at most
$(2\lfloor R\log y\rfloor+1)^4$ sites. A union bound over witnesses
therefore gives, for the discarded set $\mathcal D$,
\begin{equation}\label{eq:seed-discarded-count}
 t|\mathcal D|\le C_R\sum_{j\in\mathcal S}r_j.
\end{equation}
There is no factor $L^4$ in this count.

For the shortened path field $u_i$, use the nonnegative finite-range
coefficient matrix $k_{ij}$, whose row and column sums are at most
$C$ and which satisfies $\|u_i\|\le\sum_j k_{ij}r_j$.
Locality and the sector conditions imply $r_j\le t$ for a contributing
$j$ outside $\mathcal S$. Thus
\[
 \|u_i\|\le Ct+\sum_{j\in\mathcal S}k_{ij}r_j,
 \qquad
 \sum_{i\in\mathcal D}\|u_i\|
       \le C(C_R+1)\sum_{j\in\mathcal S}r_j.
\]
The second inequality uses \eqref{eq:seed-discarded-count} for the
constant part and the column bound for the large part. Since
$H_i\ge-\epsilon t\|u_i\|$ on discarded rows, choosing $L$ so large
that $\epsilon C(C_R+1)\le\lambda_f/4$ and combining with
\eqref{eq:seed-direct-income} proves
\begin{equation}\label{eq:seed-fine-income}
 F_A:=\sum_i(D_i+H_i)
 \ge \frac{\lambda_f}{2}t\sum_{i\in\mathcal S}r_i
             +d_f t^2|\mathcal F\setminus\mathcal S|,
 \qquad d_f=a_f/128>0.
\end{equation}
The geometric row and column estimates and the guide bound used here
are explicit inputs; the discarded-row argument shows precisely how
they enter the energy estimate.

For predictor defects the argument is simpler. Let $p>0$ be the
predictor penalty and $z_i$ its bond chord. A predictor energy row is
at least $(\sqrt p\,z_i-\|w_i\|)^2/2$, with $\|w_i\|\le Ct$.
On each listed predictor defect $i\in\mathcal I$, the sector condition
is $z_i\ge Wt$. Choose $W$ after $L$ so that $C\le\sqrt p\,W/2$.
The unlisted rows are nonnegative, and hence
\begin{equation}\label{eq:seed-predictor-income}
 I_A\ge \frac{p}{8}\sum_{i\in\mathcal I}z_i^2.
\end{equation}

\subsection{The reserve for very large defects}
Let $T\ge t$ be a larger threshold and put
\[
 \mathcal F_T=\{i\in\mathcal F:r_i>T\},\quad
 \mathcal I_T=\{i\in\mathcal I:z_i>T\},\quad
 d_A=|\mathcal F_T|+|\mathcal I_T|.
\]
The use of $T$ here does not change the seed set. Let $m\ge0$ satisfy
\[
 m\le\lambda_f/8,\qquad m\le d_f,\qquad m\le p/32.
\]
If $c_F\le\lambda_f/4$, then \eqref{eq:seed-fine-income} implies
\begin{equation}\label{eq:seed-fine-reserve}
 F_A\ge mt^2|\mathcal F|+mtT|\mathcal F_T|
                     +c_Ft\sum_{i\in\mathcal S}r_i.
\end{equation}
To check this, both $t|\mathcal S|$ and $T|\mathcal F_T|$ are at
most $\sum_{\mathcal S}r_i$. Their two charges and the $c_F$ charge
use at most $(2m+c_F)t\sum_{\mathcal S}r_i\le
(\lambda_f/2)t\sum_{\mathcal S}r_i$. The remaining fine seeds are
paid by $d_ft^2$.

Similarly suppose $W\ge1$ and $C_1\le pW/16$. Then
\eqref{eq:seed-predictor-income} implies
\begin{equation}\label{eq:seed-predictor-reserve}
 I_A\ge mt^2|\mathcal I|+mtT|\mathcal I_T|
                         +C_1t\sum_{i\in\mathcal I}z_i.
\end{equation}
Indeed $z_i\ge Wt\ge t$ on every listed bond, so
$t^2|\mathcal I|\le\sum z_i^2$ and
$tT|\mathcal I_T|\le\sum z_i^2$. Also
$C_1t z_i\le(p/16)z_i^2$. The two first charges use at most
$2m\sum z_i^2\le(p/16)\sum z_i^2$, leaving the required linear
charge within the total $p/8$.

\subsection{Coarse flags and the affected region}
Let $\mathcal O$ be the coarse plaquette rows in the $2B$ neighbourhood
of all seeds, and let $\mathcal H\subseteq\mathcal O$ be the rows
affected by a very large defect. Split the coarse plaquettes into
\[
 \mathcal C_0=\{r_i\le s\},\quad
 \mathcal C_1=\{s<r_i\le\delta\},\quad
 \mathcal C_2=\{r_i>\delta\},
 \qquad i\in\mathcal O.
\]
The geometry supplies a constant $K_0$, fixed before $L,W$, such that
\begin{equation}\label{eq:seed-geometric-counts}
 \mathcal C_2\subseteq\mathcal H,\quad
 |\mathcal C_2|\le K_0d_A,\quad
 n_g\le4(|\mathcal C_1|+K_0d_A),\quad
 |\mathcal O|\le K_0B^4 n,
\end{equation}
where $n_g$ is the number of coarse flag seeds and
$n=|\mathcal F|+|\mathcal I|+n_g$ is the total number of seeds.
These are the precise geometric counting inputs used below.
For $t\le2s$ and $m\le a_c/64$, the middle bound in
\eqref{eq:seed-three-ranges} gives at least $4mt^2$ on each row of
$\mathcal C_1$. Summing all three ranges yields
\begin{equation}\label{eq:seed-coarse-reserve}
 C_A\ge4mt^2|\mathcal C_1|
       -\lambda_c s\sum_{i\in\mathcal C_2}r_i
       -B_0s^4|\mathcal O|.
\end{equation}

The remaining two comparison rows, denoted together by $L_A$, are
split over $\mathcal H$ and its complement. Their local estimates
have the form
\begin{align}
 L_{\mathcal H}&\ge\lambda_c s\sum_{i\in\mathcal C_2}r_i
 -\frac{C_1}{L^2}t\sum_{i\in\mathcal S}r_i
 -C_1t\sum_{i\in\mathcal I}z_i-Zt^2B^4d_A,
 \label{eq:seed-hit-input}\\
 L_{\mathcal O\setminus\mathcal H}&\ge-\eta t^2 n.
 \label{eq:seed-unhit-input}
\end{align}
The constant $C_1$ is fixed before $L,W$, whereas $Z$ may depend on
both. The second estimate is available for every fixed $\eta>0$ by
taking $x$ sufficiently large after the other parameters. The row
estimates producing these inequalities remain inputs here. The
positive normalization row, denoted $N_A$, satisfies $N_A\ge0$.
The total assigned energy is exactly
$F_A+I_A+C_A+N_A+L_{\mathcal H}+L_{\mathcal O\setminus\mathcal H}$.

Choose $L$ large enough for the fine estimate and
$C_1/L^2\le\lambda_f/4$, and $W$ large enough for the predictor
estimate, $W\ge1$ and $C_1\le pW/16$. Choose
\[
 0<m\le\min\{\lambda_f/8,d_f,p/32,a_c/64\},\qquad \eta=m/4.
\]
Suppose finally that
\begin{equation}\label{eq:seed-final-smallness}
 (4mK_0+ZB^4)t^2\le mtT,
 \qquad B_0s^4K_0B^4\le(m/4)t^2.
\end{equation}
Adding \eqref{eq:seed-fine-reserve}--\eqref{eq:seed-unhit-input}
cancels the three explicit linear expenses. The first inequality in
\eqref{eq:seed-final-smallness} leaves $4mK_0t^2d_A$ from the
very-large-defect reserve, and the second pays the small coarse rows.
Thus the assigned energy is at least
\[
 mt^2(|\mathcal F|+|\mathcal I|)
 +4mt^2(|\mathcal C_1|+K_0d_A)-\tfrac m4t^2n-\tfrac m4t^2n
 \ge\tfrac m2t^2n.
\]
The last step is the flag count in \eqref{eq:seed-geometric-counts}.
This is the required positive gain per seed, with $\mu=m/2$.

\subsection{The order of choices and orientation average}
For the actual thresholds put
\[
 t=\frac{P(h)}{\sqrt h},\qquad
 s=\frac{P(x)}{\sqrt x},\qquad
 T=h^{-49/100},\qquad h=x+c_L.
\]
As $x\to\infty$ at fixed $L$,
\[
 s\to0,\quad t/s\to1,\quad
 B^r\,t/T\to0,\quad B^r s\to0
 \quad\text{for every fixed }r\ge0.
\]
For the ratio $t/s$, subtract the logarithms:
$\log t=(\log\log h)^2-\frac12\log h$ and the corresponding formula
at $x$. Their difference tends to zero under a fixed shift, by the
mean value theorem, since the derivative of
$(\log\log x)^2-\frac12\log x$ tends to zero. The other statements
follow from $B$ and $P$ growing slower than every positive power,
using $-1/2+49/100=-1/100$ for $t/T$.
It follows that $0<s\le\delta$, $t\le2s$, $t\le T$, and both
inequalities \eqref{eq:seed-final-smallness} hold eventually. For the
second, divide by $t^2$ and use $s/t\to1$ and $B^4s^2\to0$.
All choices have therefore been made in the order
\[
 \text{fixed group and row data}\ ;\ L\text{ large}\ ;\
 W\text{ large}\ ;\ \mu>0\ ;\ x\text{ large},
\]
with no dependence on the lattice period in the last threshold.

The packet construction uses an average of the assigned energies over
signed coordinate transformations. Here the block side is odd, so
the parent map commutes with the signed coordinate action. We use
the patch's signed covariance laws for the predictor and the local
likelihoods. Turn the tree, seed set and tag array together with the fields. Tree edges remain equal to the group
identity; the predictor and likelihood transform covariantly; distances,
chords and the number of seeds are preserved. The local sector
conditions, including the strict positivity of the retained site
likelihoods, therefore hold in every transformed channel. Applying the
same energy bound there and averaging with normalized weights yields
$Q_A^{\rm core}\ge\mu t^2n$. This supplies the core input to
\eqref{eq:packet-assigned-energy} once the geometric, guide and
comparison inputs specified in this section have been established.
The average is of energies on transformed copies of the same data;
there is no independent choice of trees or tags between channels.
